\documentclass[12pt,a4paper]{article}
\usepackage[utf8]{inputenc}
\usepackage{amsmath, amssymb, amsfonts}
\usepackage{geometry}
\geometry{margin=2.5cm}
\usepackage{hyperref}
\usepackage{pgfplots}
\usepackage{fancyhdr}
\usepackage{listing}
\usepackage{listings}
\usepackage{titlesec}
\titleformat{\section}
  {\normalfont\Large\bfseries\centering}  % Format: font size, bold, centered
  {\thesection}{1em}{}
\pagestyle{fancy}
\pgfplotsset{compat=1.18}
\fancyhead{}
\fancyfoot{}
\title{Architetture e Sistemi Operativi}
\author{}
\date{}
\fancyhead[L]{Gabriele Perna}          % Left side
\fancyhead[C]{Architetture e Sistemi Operativi}     % Center
\fancyhead[R]{2026/2027}             % Right side
\fancyfoot[C]{\thepage}           % Page number in the center of footer
\let\oldsection\section
\renewcommand{\section}{\clearpage\oldsection}

\begin{document}

\maketitle
\tableofcontents

\section{Introduzione}
In questo corso andremo a:
\begin{itemize}
    \item Capire come funziona un elaboratore.
    \item Capire quali tipi di istruzioni è in grado di eseguire un elaboratore nativamente.
    \item Come lo si può dotare di un sistema operativo.
\end{itemize}
Il corso utilizza Windows e Linux come esempi di sistemi operativi.
\paragraph{Codice Teams:} 6uk5ntg
\paragraph{Email Docente:} \href{marco.delutto@unipi.}{marco.delutto@unipi.}

\section{Processore, memoria e istruzioni}
\subsection{Processore}
Il \textbf{processore} è in grado di fare operazioni piccole.\\
Il processore contiene i \textbf{registri}, memorie piccolissime e velocissime utilizzate per memorizzare dati temporanei legati alle istruzioni da eseguire.
Ecco alcuni registri speciali:
\begin{itemize}
    \item \textbf{Program Counter (PC)}: un registro speciale che conserva lo stato di esecuzione memorizzando l'indirizzo della prossima istruzione da eseguire.
    \item \textbf{Instruction Register (IR)}: contiene l'istruzione prelevata dalla memoria che deve essere eseguita.
    \item \textbf{Stack Pointer (SP)}: contiene la posizione attuale dello stack.
\end{itemize}
\subsection{Memoria}
Il processore è unito ad una \textbf{memoria}, costituita da \textbf{celle}, è vista come un vettore:
\[
M[i] = k
\]
Ogni cella ha un valore e la sua capacità è 1 parola (32 o 64 bit).\\
Una volta che il processore viene acceso, esso esegue:
\begin{lstlisting}
    while(true):
        instr = getInstructionFrom(M)
        decoded = decode(instr)
        execute(decoded)
        determineResults()
        ...
\end{lstlisting}
Il processore quindi:
\begin{itemize}
    \item Estrae istruzioni dalla memoria
    \item Le decodifica
    \item Le esegue
    \item Calcola i risultati
\end{itemize}

\subsection{Tipi di istruzioni}
Le istruzioni possono essere:
\begin{itemize}
    \item Operative
    \begin{itemize}
        \item Aritmetiche
        \begin{itemize}
            \item ADD
            \item SUB
            \item MUL
            \item DIV
        \end{itemize}
        \item Logiche
        \begin{itemize}
            \item AND
            \item OR
            \item NOT
        \end{itemize}
        \item Shift / rotazione
        \begin{itemize}
            \item SHIFT
            \item ROT
        \end{itemize}
        \item \dots
    \end{itemize}
    \item Di Memoria
    \begin{itemize}
        \item LOAD
        \begin{itemize}
            \item Restituisce un valore da un indirizzo:\\
            \begin{lstlisting}
                LOAD val,addr
            \end{lstlisting}
        \end{itemize}
        \item STORE
        \begin{itemize}
            \item Salva un valore in un indirizzo:\\
            \begin{lstlisting}
                STORE val,addr
            \end{lstlisting}
        \end{itemize}
    \end{itemize}
    \item di Salto
    \begin{itemize}
        \item "Saltano" pezzi di codice, volendo anche secondo istruzioni, aggiornando il PC alla corrispondente istruzione.\\
        \begin{lstlisting}
            inizio:
            \dots
            jmp fine:
            -skipped code-
            fine:
        \end{lstlisting}
    \end{itemize}
\end{itemize}

\section{Livelli di astrazione dei calcolatori}
Per capire come opera un computer osserviamo dei livelli di astrazione, sviluppati a "torre" nel seguente modo:
\begin{center}
    applicazioni\\
    sistema operativo\\
    architettura (set istruzioni)\\
    $\mu$ - architettura (processore)\\
    logica (componenti)\\
    circuiti digitali (0/1, 5V/0V)\\
    circuiti analogici\\
    dispositivi\\
    fisica
\end{center}
Non andremo a vedere niente delle caratteristiche hardware degli ultimi 3 livelli, non ci interessa.\\
Per capire:
\begin{itemize}
    \item Noi partiamo col realizzare dei circuiti digitali che servono per creare dei componenti
    \item Tali componenti ci servono per realizzare un processore
    \item Al processore dovremo dare in pasto delle istruzioni
    \item Le istruzioni servono al sistema operativo e alle applicazioni
\end{itemize}
Tra i livelli di "architettura" e "sistema operativo" si trova un \textbf{interfaccia}, genericamente attribuita al livello \textbf{i}. Questo significa che:
\begin{itemize}
    \item Il livello del sistema operativo corrisponde al livello \textbf{i} (si trova al livello dell'interfaccia)
    \item Il livello dell'architettura corrisponde al livello \textbf{i - 1}
\end{itemize}
Al livello \textbf{i + 1}, tra sistema operativo e applicazioni, si trovano le \textbf{chiamate al sistema operativo}, chiamate \textbf{API}.\\\\
Ogni livello conosce solo il \textbf{livello inferiore} a se, questo è per mantenere i livelli superiori attraverso diversi livelli inferiori; per esempio un telefono Samsung e un Motorola sono realizzati in maniera diversa, ma entrambi utilizzano Android, questo è possibile proprio grazie al fatto che i livelli non si concentrano sui livelli soprastanti.\\
Un altro esempio: non si possono portare app iOS su Android.\\
Spesso esistono dei traduttori che permettono di convertire istruzioni (x64,x86) per permettere compatibilità con sistemi operativi e processori differenti.
\subsection{Principi importanti}
I componenti che svolgono le istruzioni seguono questi concetti:
\begin{itemize}
    \item \textbf{Gerarchia}
    \item \textbf{Modularità}: presenza di moduli per la specializzazione e composizione del lavoro.
    \item \textbf{Regolarità}
\end{itemize}
Si usano moduli più piccoli per svolgere azioni più grandi quando combinati.

\section{Circuiti Digitali}
\subsection{Sistema binario}
Noi utilizziamo il sistema decimale, ovvero un sistema \textbf{posizionale}, dove ogni cifra possiede un valore diverso a seconda della sua posizione.\\
Anche il \textbf{sistema binario} è un sistema posizionale, ma le cifre sono solo 2 (0,1).\\
\begin{itemize}
    \item \textbf{Esempio decimale:}
    \[
    13_{10} = (1 \cdot 10^1) + (3 \cdot 10^0)
    \]
    \item \textbf{Esempio binario:}
    \[
    10_2 = (1 \cdot 2^1) + (0 \cdot 2^0) = 2_{10}
    \]
\end{itemize}
Nei computer si utilizza il sistema binario per via della sua semplicità e della sua compatibilità con i moduli logici, per esempio una moltiplicazione è più facile da svolgere in binario tramite un modulo \textbf{AND}, invece che in decimale.\\\\
Per rappresentare i numeri negativi in binario ci sono due modi:
\begin{itemize}
    \item \textbf{Modulo e segno:} si riserva il \textbf{MSB} (Most Significant Bit, più a sinistra) per il segno (1 se negativo, 0 se positivo).
    \item \textbf{Complemento a 2:} nato perche complicate le operazioni con modulo e segno (if else); si ha un algoritmo per ottenere il binario di un numero negativo:
    \begin{enumerate}
        \item Si trasforma il numero in binario.
        \item Si fa il NOT di ogni bit
        \item si aggiunge "+1"
    \end{enumerate}
\end{itemize}
\end{document}